% Propietario conservado: /Users/ruben/Documents/ChatGPT/jueces y controles/REVISION_ARGUMENTAL_APERTURAS_CIERRES_20260915/II/manuscript_es/sections/18_comparacion_cuantitativa_II.tex
\section{Évaluation des constantes et comparaison métrologique}
\label{ii:sec:ii-contraste-cuantitativo}

Les constructions précédentes permettent de séparer trois opérations :
produire une grandeur, publier sa coordonnée dans une carte et comparer cette
coordonnée à une référence. L’état APP--TRIT--TPK conserve les feuillets,
l’orientation, le quotient, la retenue et la mémoire avant ces lectures. Sa
prolongation nonadique publie les coordonnées corrélées
\(\pi,\varphi,e,\alpha\) ; le lecteur déterminantal et le retour régional
produisent les sections d’action ; l’incidence marquée \(90/120\) et
les canaux conjugués produisent la fonctionnelle de Barbero ; le transport
angulaire publie ensuite les coordonnées CKM. Aucune entrée des tables
externes n’intervient dans le choix de ces opérations.

Sont réunies ici les évaluations conservées dans les documents propriétaires et leurs
comparateurs. Les chiffres des sorties internes expriment la précision de
publication numérique, et non une incertitude expérimentale. Pour une valeur
\(x\) et une référence non nulle \(x_0\), toutes deux dans la même carte, on utilise
\[
 \Delta x=x-x_0,\qquad r_x=\frac{x-x_0}{x_0}.
\]
Lorsque la référence possède des incertitudes marginales asymétriques
\(u_-,u_+\), nous définissons en outre le résidu descriptif
\begin{equation}
 z_x=\begin{cases}
 \Delta x/u_+,&\Delta x\geq0,\\
 \Delta x/u_-,&\Delta x<0.
 \end{cases}
 \label{ii:eq:ii-residuo-marginal}
\end{equation}
Cette normalisation ne présuppose pas l’indépendance entre les lignes et ne constitue pas
à elle seule un test statistique conjoint.

\subsection{Constante de Planck et constante réduite : deux orientations}

Dans la carte \(\mathcal U_S\mapsto\mathrm{J\,s}\), les sections déjà
construites sont
\begin{align}
 \hbar_{\rm pre}&=H_5\,10^{-34}\,\mathrm{J\,s},&
 H_5&=1.0545718176461544696\ldots,\\
 \hbar_{\rm ret}&=\eta_{\rm ret}\,10^{-34}\,\mathrm{J\,s},&
 \eta_{\rm ret}&=1.0545718169150390611\ldots,\\
 h_a&=2\pi\hbar_a,&a&\in\{\mathrm{pre},\mathrm{ret}\}.
 \label{ii:eq:ii-cuatro-coordenadas-planck}
\end{align}
La décade provient de \(\Sigma_H=\varphi\alpha^{16}\) ; la différence entre
les deux sections provient du retour régional. Multiplier par \(2\pi\)
fait passer de l’action par radian à l’action par tour, sans introduire une autre
orientation ni altérer le résidu relatif.

La référence du SI est
\[
 h_{\rm SI}=6.62607015\times10^{-34}\,\mathrm{J\,s},\qquad
 \hbar_{\rm SI}=\frac{h_{\rm SI}}{2\pi}
 =1.0545718176461563912\ldots\times10^{-34}\,\mathrm{J\,s}.
\]
Les deux expressions sont exactes. Le développement infini de \(\hbar_{\rm SI}\)
n’est pas une incertitude de mesure. La fixation de la valeur numérique de
\(h\) appartient à la définition du SI en vigueur depuis 2019
\cite{ii:bipm2018} ; la comparaison s’effectue après la construction HMT.

\begin{table}[htbp]
\centering\small
\setlength{\tabcolsep}{5pt}
\begin{tabular}{@{}llr@{}}
\toprule
Magnitud & Coordonnée en \(10^{-34}\,\mathrm{J\,s}\) & \(r_x\) par rapport au SI\\
\midrule
\(\hbar_{\rm pre}\) & \(1.0545718176461544696\) & \(-1.822221\times10^{-15}\)\\
\(\hbar_{\rm ret}\) & \(1.0545718169150390611\) & \(-6.932836\times10^{-10}\)\\
\(h_{\rm pre}\) & \(6.62607014999998793\) & \(-1.822221\times10^{-15}\)\\
\(h_{\rm ret}\) & \(6.62607014540625433\) & \(-6.932836\times10^{-10}\)\\
\bottomrule
\end{tabular}
\caption{Action avant et après le retour. Les derniers chiffres sont des
évaluations des coordonnées publiées ; les résidus ne sont pas des erreurs
expérimentales de \(h\) ni de \(\hbar\).}
\label{ii:tab:ii-planck-comparacion}
\end{table}

\begin{proposition}[Invariance de la comparaison sous le tour complet]
Dans une carte commune,
\[
 \frac{h_a-h_{\rm SI}}{h_{\rm SI}}
 =\frac{\hbar_a-\hbar_{\rm SI}}{\hbar_{\rm SI}}.
\]
\end{proposition}
\begin{proof}
Les deux différences du premier membre contiennent le facteur commun
\(2\pi\), qui s’annule avec celui du dénominateur.
\end{proof}
La séparation pre/ret est quantitativement pertinente : remplacer l’une par
l’autre change la comparaison de plusieurs ordres de grandeur. Le correcteur
régional doit rester visible et ne pas être confondu avec un arrondi numérique.

\subsection{Barbero--Immirzi : valeur interne et comparaison entre constructions}

L’évaluation de la fonctionnelle se reconstruit avec ses trois contributions :
\begin{align*}
 \Delta S_{90}&=2.95169439297457621139847987861\times10^{-4},\\
 \Delta S_{120}&=1.96835112502951720093738706217\times10^{-5},\\
 \pi AC^*/180&=0.270485322934140078465586458790\ldots,\\
 \gamma_{\rm HMT}&=\pi AC^*/180+12\Delta S_{90}-\Delta S_{120}\\
 &=0.27400767269445927474725526077398041940\ldots.
\end{align*}
Cette composition conserve le poids orienté \(12:-1\), antérieur à son
évaluation décimale. Ses termes publiés reconstruisent la coordonnée
finale avec un résidu absolu inférieur à \(3\times10^{-31}\).

Comme comparaison bibliographique, on considère le dénombrement de Ghosh et Mitra,
non une mesure du paramètre. Leur équation (7), écrite avec
\(\lambda=2\pi\gamma_{\rm GM}\), est
\begin{equation}
 \sum_{j\in\{1/2,1,3/2,\ldots\}}(2j+1)
 e^{-2\pi\gamma_{\rm GM}\sqrt{j(j+1)}}=1.
 \label{ii:eq:ii-gm-comparador}
\end{equation}
La référence est A.~Ghosh et P.~Mitra,
\emph{An improved estimate of black hole entropy in the quantum geometry
approach}, \emph{Physics Letters B} \textbf{616} (2005), 114--117,
\href{https://arxiv.org/abs/gr-qc/0411035}{arXiv:gr-qc/0411035},
équations (7)--(8). La racine est évaluée ici pour cette équation précise ;
on n’identifie pas le dénombrement des ponctures au générateur HMT.

\begin{table}[htbp]
\centering\small
\begin{tabular}{@{}lll@{}}
\toprule
Construction & Valeur adimensionnelle & Nature\\
\midrule
Fonctionnelle \(\Gamma_{6\to8}\) & \(0.2740076726944593\) & Incidence et retour HMT\\
Racine de \eqref{ii:eq:ii-gm-comparador} & \(0.2740668582328300\) & Dénombrement bibliographique déclaré\\
\midrule
\(\gamma_{\rm HMT}-\gamma_{\rm GM}\) & \(-5.918553837\times10^{-5}\) & Différence entre constructions\\
\(\gamma_{\rm HMT}/\gamma_{\rm GM}-1\) & \(-2.159529202\times10^{-4}\) & Diferencia relativa\\
\bottomrule
\end{tabular}
\caption{Comparaison de Barbero--Immirzi dans l’orientation positive. Aucune
incertitude expérimentale ni signification statistique n’est assignée à ce tableau.}
\label{ii:tab:ii-barbero-comparacion}
\end{table}

L’évaluation de la racine externe est reproductible sans tronquer
arbitrairement sa queue. Avec \(n=2j\), \(N=128\) et
\(r=e^{-0.27\pi}\), pour \(\gamma\geq0.27\),
\[
 \sum_{n>N}(n+1)e^{-\pi\gamma\sqrt{n(n+2)}}
 \leq \frac{r^{N+1}[(N+2)-(N+1)r]}{(1-r)^2}<10^{-43}.
\]
Chaque terme décroît strictement avec \(\gamma\). Les extrémités
\(0.27\) et \(0.28\) encadrent la racine ; la dichotomie et la borne de queue
contrôlent les décimales du tableau. Il s’agit d’une vérification du
comparateur bibliographique, non d’une voie pour sélectionner \(\gamma_{\rm HMT}\).

\subsection{Angles CKM, amplitudes et violation CP}

Le lecteur angulaire déjà construit publie, en degrés,
\begin{equation}
 (\theta_{12},\theta_{23},\theta_{13},\delta)
 =(13.003313589509,\;2.398056157332,\;
 0.213977382244,\;65.676173123554)^\circ.
 \label{ii:eq:ii-ckm-decimales}
\end{equation}
Ces valeurs proviennent de \(M_{\rm CKM}(A,C^*/6,\Delta)^T\) et
\(\delta=9A\). La conversion en radians s’effectue une seule fois avant
d’évaluer les sinus et les cosinus. La paramétrisation unitaire de la section
angulaire donne
\[
 |V_{\rm CKM}|\simeq
 \begin{pmatrix}
 0.974350259&0.225005836&0.003734601\\
 0.224873110&0.973489279&0.041841465\\
 0.008582400&0.041121745&0.999117283
 \end{pmatrix},
\]
et le quartet invariant conserve la phase :
\[
 J=c_{12}c_{23}c_{13}^{\,2}s_{12}s_{23}s_{13}\sin\delta
 =3.1189723326632974\times10^{-5}\quad
 \text{dans l'évaluation du propriétaire.}
\]
Cette évaluation utilise \eqref{ii:eq:ii-jarlskog} avec les coordonnées de
\eqref{ii:eq:ii-ckm-decimales}. Les neuf modules imprimés sont arrondis ; l’unitarité appartient
à la matrice complexe avant l’arrondi, non à une matrice formée
uniquement de modules.

Le comparateur est l’ajustement à trois générations du
\href{https://pdg.lbl.gov/2025/reviews/rpp2025-rev-ckm-matrix.pdf}
{Particle Data Group, mise à jour de 2025, \emph{CKM Quark-Mixing Matrix}},
les équations (12.27)--(12.28) et la valeur voisine de \(J\). Pour éviter
d’altérer leurs incertitudes par une transformation non linéaire, le tableau
conserve les observables publiées par PDG : \(s_{ij}=\sin\theta_{ij}\),
\(\delta\) en radians et \(J\).

\begin{table}[htbp]
\centering\small
\setlength{\tabcolsep}{4pt}
\begin{tabular}{@{}lrrr@{}}
\toprule
Observable & Évaluation HMT & PDG 2025 & \(z_x\)\\
\midrule
\(s_{12}\) & \(0.225007404757\) & \(0.22501\pm0.00068\) & \(-0.003817\)\\
\(s_{23}\) & \(0.041841757010\) & \(0.04183^{+0.00079}_{-0.00069}\) & \(0.014882\)\\
\(s_{13}\) & \(0.003734601164\) & \(0.003732^{+0.000090}_{-0.000085}\) & \(0.028902\)\\
\(\delta\) (rad) & \(1.146265461116\) & \(1.147\pm0.026\) & \(-0.028251\)\\
\(10^5J\) & \(3.118972333\) & \(3.12^{+0.13}_{-0.12}\) & \(-0.008564\)\\
\bottomrule
\end{tabular}
\caption{Comparaison marginale dans la paramétrisation et l’édition déclarées.
Le facteur \(10^5\) multiplie conjointement la valeur et l’incertitude de la
dernière ligne ; il ne modifie pas son résidu normalisé.}
\label{ii:tab:ii-ckm-comparacion}
\end{table}

Les résidus sont calculés avec les coordonnées avant l’arrondi
du tableau. L’ajustement PDG combine données et contraintes d’unitarité ;
ses lignes ne sont pas des mesures indépendantes. On ne somme donc pas
\(\sum z_x^2\) et on n’attribue pas de nombre de degrés de liberté à ce tableau.
Le résultat est une comparaison marginale reproductible avec cette édition,
et non une mesure de probabilité du modèle. Les écarts marginaux
historiques conservés dans le document propriétaire ne sont pas reportés dans ce tableau :
l’extrait n’identifie pas l’édition ni la covariance de cette comparaison.

Le secteur PMNS conserve le transport \(U=F_\ell^\dagger F_\nu\), les
domaines spectraux et les données de phase développés dans la section
angulaire. Une ligne numérique exige de préciser les repères du registre
utilisé et la réalisation comparée. Ce tableau n’attribue pas de chiffres CKM
au lecteur leptonique et ne publie pas de nouvelle matrice PMNS.

\subsection{Boltzmann et gravitation universelle : grandeur, carte et référence}

La construction thermique déjà retrouvée détermine le transducteur positif
\(k_B:\mathcal L_\Theta\to\mathcal L_E\) et la relation
\begin{equation}
 k_B\Theta_{{\rm clk},a}\log3
 =\frac{h_a}{108t_0}=\frac{\hbar_a}{t_P},
 \qquad t_P=\frac{54}{\pi}t_0.
 \label{ii:eq:ii-kb-comparacion}
\end{equation}
L’information locale, l’action et la période apparaissent avant la carte
kelvin--joule. Pour des bases positives \(u_E,u_\Theta\), si
\(k_B(u_\Theta)=b\,u_E\),
\(\Theta_{{\rm clk},a}=\theta_a u_\Theta\) et
\(E_{P,a}=\epsilon_a u_E\), l’équation équivaut à
\(b\theta_a=\epsilon_a/\log3\). Un changement
\(u'_E=s_Eu_E\), \(u'_\Theta=s_\Theta u_\Theta\) produit
\[
 b'=\frac{s_\Theta}{s_E}b,\qquad
 \theta'_a=\frac{\theta_a}{s_\Theta},\qquad
 \epsilon'_a=\frac{\epsilon_a}{s_E}.
\]
Le produit énergétique est ainsi conservé et l’opération
supplémentaire qui sépare les coordonnées individuelles devient visible. Pour le transducteur
de la sous-section~\ref{ii:sec:ii-kb-aut-desarrollo}, la carte kelvin--joule
utilisée ici exprime la coordonnée SI
\[
 k_B^{\rm SI}=1.380649\times10^{-23}\,\mathrm{J\,K^{-1}},
\]
valeur exacte de la définition du SI, et non estimation assortie d’une erreur de mesure.
La fonction du transducteur et sa relation avec l’action ont un contenu
propre ; citer sa coordonnée SI ne constitue pas une vérification numérique
indépendante d’un choix interne des deux bases.

En gravité, la réalisation circulaire fixe
\begin{equation}
 G_a=\vartheta_{108}^{\circ}\frac{c_{\rm int}^{5}t_0^2}{\hbar_a},
 \qquad
 \vartheta_{108}^{\circ}=(54/\pi)^2
 =295.4525715010569\ldots.
 \label{ii:eq:ii-g-comparacion}
\end{equation}
Le sélecteur et la réciprocité \(G_a\hbar_a=
\vartheta_{108}^{\circ}c_{\rm int}^{5}t_0^2\) sont des grandeurs internes de
la réalisation déclarée. Une évaluation dimensionnelle de \(G_a\) doit
conserver, outre l’orientation \(a\), les cartes de temps,
de vitesse et d’action. Dans la carte dimensionnelle utilisée ici, la valeur évaluée est
\(6.67430\times10^{-11}\,\mathrm{m^3\,kg^{-1}\,s^{-2}}\).
Le tableau distingue cette coordonnée de la
valeur recommandée avec son incertitude.

\begin{table}[htbp]
\centering\small
\begin{tabularx}{\linewidth}{@{}lXX@{}}
\toprule
Objeto & Sortie ou relation interne & Referencia dimensional\\
\midrule
\(k_B\) & Transducteur positif ; \(k_B\Theta_{{\rm clk},a}=E_{P,a}/\log3\).
& \(1.380649\times10^{-23}\,\mathrm{J/K}\), exact dans le SI.\\[4pt]
\(G_a\) & Sélecteur \((54/\pi)^2\) ; \(G_a=\vartheta_{108}^{\circ}c_{\rm int}^5t_0^2/\hbar_a\).
& \(6.67430(15)\times10^{-11}\,\mathrm{m^3/(kg\,s^2)}\), CODATA 2022.\\
\bottomrule
\end{tabularx}
\caption{Normalisation interne et référence de carte. L’incertitude-type
relative de la référence de \(G\) est d’environ
\(2.24743\times10^{-5}\) ; \(k_B\) n’a pas d’incertitude définitionnelle.
On ne calcule pas de résidu prédictif à partir d’une coordonnée déjà
identifiée comme référence de carte.}
\label{ii:tab:ii-kb-g-cartas}
\end{table}

Les sources externes de ce dernier tableau sont les
\href{https://www.bipm.org/en/measurement-units/si-defining-constants}
{constantes définissant le SI du BIPM} et le
\href{https://physics.nist.gov/cuu/Constants/Table/allascii.txt}
{tableau CODATA 2022 du NIST} \cite{ii:codata2022}.
L’incertitude expérimentale de \(G\) ne se transfère pas à \(h\),
\(\hbar\) ni à \(k_B\). Une égalité de chiffres dans une carte
ne supprime pas non plus la nécessité d’identifier le lecteur qui produit la grandeur.

\subsection{Reproduction des évaluations}

Le programme \texttt{verificar\_comparaciones\_II.py}, inclus avec les
sources, reçoit les coordonnées internes publiées dans des fichiers distincts
des données externes. Il vérifie la composition de la fonctionnelle de Barbero,
la conversion entre \(h\) et \(\hbar\), la racine bibliographique et sa queue,
les modules CKM, le quadruplet \(J\) et chaque résidu imprimé. Il utilise
une arithmétique décimale à 75 chiffres et des fonctions élémentaires évaluées après
le gel des sorties internes. Les instructions d’exécution
et les versions des programmes sont incluses dans le paquet reproductible.
La précision de travail n’augmente pas le nombre de chiffres disponibles dans le
document propriétaire : seuls sont publiés les chiffres stables sous son arrondi.
Le contrôle enregistre les fichiers utilisés et leurs empreintes, afin
qu’un changement d’évaluation ou d’édition externe reste visible.

% NUMCHECK hbar_SI_coefficient=1.0545718176461563912 tol=1e-19
% NUMCHECK h_pre_coefficient=6.62607014999998793 tol=6e-18
% NUMCHECK h_ret_coefficient=6.62607014540625433 tol=6e-18
% NUMCHECK relative_pre=-1.822221e-15 tol=6e-22
% NUMCHECK relative_ret=-6.932836e-10 tol=6e-17
% NUMCHECK gamma_GM=0.2740668582328300 tol=6e-17
% NUMCHECK gamma_difference=-5.918553837e-5 tol=6e-15
% NUMCHECK gamma_relative=-2.159529202e-4 tol=6e-14
% NUMCHECK ckm.s12.value=0.225007404757 tol=6e-13
% NUMCHECK ckm.s23.value=0.041841757010 tol=6e-13
% NUMCHECK ckm.s13.value=0.003734601164 tol=6e-13
% NUMCHECK ckm.delta_rad.value=1.146265461116 tol=6e-13
% NUMCHECK ckm.s12.marginal_residual=-0.003817 tol=6e-7
% NUMCHECK ckm.s23.marginal_residual=0.014882 tol=6e-7
% NUMCHECK ckm.s13.marginal_residual=0.028902 tol=6e-7
% NUMCHECK ckm.delta_rad.marginal_residual=-0.028251 tol=6e-7
% NUMCHECK ckm.J.marginal_residual=-0.008564 tol=6e-7
% NUMCHECK gravity_selector=295.4525715010569 tol=6e-14
% NUMCHECK G_relative_reference_uncertainty=2.24743e-5 tol=6e-11
